% =============================================================================
%  三、问题求解（四）：系统对比  +  四、结论
% =============================================================================
\subsection{新旧方法的系统对比}

\subsubsection{相同的部分}

替换方法最容易走偏的地方，是把“换方法”变成“换问题”。为保证可比性，本文
在三件事上刻意保持不变：\textbf{同一数据集}（519 名学生的 14 个题项与成
绩）、\textbf{同一评价对象}（五类学生、五级效率层级）、\textbf{同一目标}
（让个体从低效层升到高效层）。甚至连“层”的定义也沿用原文的分层思想，只是
把剔除式换成了分位带，以保证状态空间占满。因此，下文报告的差异可以归因于
方法，而不是数据或问题的变化。

\subsubsection{不同的部分}

两种范式在六个维度上根本不同（表 \ref{tab:compare}）。

\begin{table}[htbp]
\centering
\tabfont
\caption{context-dependent DEA 与随机过程范式的对比}
\label{tab:compare}
\begin{tabularx}{\textwidth}{@{}p{2.9cm}p{4.2cm}X@{}}
\toprule
维度 & 原文：CDEA & 本文：随机过程 \\
\midrule
基本对象 & 截面上的个体与参照集 &
  个体在时间上的状态轨道 \\
核心量 & 效率 $\theta$、层号 $S$、三准则 &
  强度 $\rho$、转移矩阵 $\bP(t)$、首达时间 $m_k$ \\
时间维度 & 无（“一学期”只是背景） &
  显式（周为单位，可算任意时长） \\
输出形态 & 排名、层号、四步参考路径 &
  跃层概率、期望时间、最优动作组合 \\
对“投入”的态度 & 投入越多越接近前沿（若产出同增） &
  投入增加会稀释效率，方向由 $\delta$ 的符号决定 \\
数据要求 & 单期截面即可 &
  截面可估相对强度，绝对尺度需锚点或面板 \\
主要风险 & 口径依赖（层数、$k$、权重） &
  识别性弱（无面板时强度不可全识别） \\
\bottomrule
\end{tabularx}
\end{table}

\textbf{（1）从“位置”到“过程”。}原文的核心输出是“你在第几层、离上一层
多远、下一个模仿谁”；本文的核心输出是“你多久能上一层、上行概率多大、做
什么能提高上行强度”。前者是空间性的，后者是时间性的。

\textbf{（2）从“排序”到“概率”。}原文的路径是确定性的：给定数据，路径
唯一。本文的结论是概率性的：16 周末有 64.2\% 的概率已不在第 5 层，也有
35.8\% 的概率原地不动。概率化的表述更诚实，因为它承认“同样的建议在不同
人身上会有不同结果”。

\textbf{（3）从“模仿”到“选择”。}原文的推荐规则是“向上一层最像你、最容
易跟上的人学习”；本文的推荐规则是“在预算约束下选择使上行强度最大的动作
组合”。前者是一个启发式，后者是一个优化问题。

\textbf{（4）从“单一指标”到“参数解释”。}原文的三准则（$\Omega,P,D,H$）
是复合指标，其数值难以直接解释为行为含义；本文的 $\delta_j$ 是回归系数，
可以直接读出“复习每提高一个标准差，上行—下行强度比变化多少”，从而支持
“为什么”的追问。

\textbf{（5）代价。}随机过程范式并非免费：它要求时间齐次、马尔可夫、层内
同质、协变量线性等假设，要求一个外部锚点来定标时间，且在没有面板数据时无
法识别个体层面的转移强度。这些假设在本文第三节被逐条列出，并在结论中作为
局限重申。

\subsubsection{数值层面的对话}

两个范式在同一名个体上给出的数字可以直接对照（表 \ref{tab:numbers}）。有
三处值得注意。

\begin{table}[htbp]
\centering
\tabfont
\caption{S480 个案：两种范式给出的数值对照}
\label{tab:numbers}
\begin{tabularx}{\textwidth}{@{}p{3.2cm}p{4.6cm}X@{}}
\toprule
指标 & 原文（DEA 路径） & 本文（随机过程） \\
\midrule
当前效率 & 0.6771 & 0.6771 \\
单步提升（16 周） & 25.38\% & 25.53\% \\
累计提升 & 43.97\%（4 步） & 47.69\%（路径）／43.97\%（53.7 周） \\
所需时间 & 未给出 & 14.5 周（首达最高层）／53.7 周（累计 43.97\%） \\
建议动作 & 模仿 S477，+4.1 h/周 & 兴趣拓展，+1.5 h/周 \\
建议后的跃层概率 & 未给出 & 0.781（最优）／0.617（模仿） \\
\bottomrule
\end{tabularx}
\end{table}

\textbf{第一，单步增益几乎重合。}25.38\% 与 25.53\% 相差 0.15 个百分点。
这不是因为两种方法等价，而是因为本文用前者标定了后者的时间尺度；它证明了
“一学期 $\approx$ 一次跃层”这一换算在数量上站得住，从而让后文的“周”具
有可信的物理含义。

\textbf{第二，累计增益出现分歧。}原文四步累计 43.97\%，本文路径累计
47.69\%，相差 3.7 个百分点；而本文的 OU 模型反解 43.97\% 需要 53.7 周。三
个数字放在一起说明：原文的“四步”是在层级空间中跳跃，不考虑每步之间的时
间成本；本文的“53.7 周”是把同样的增益摊到连续时间上。前者是“最好的情
形”，后者是“平均的情形”。

\textbf{第三，建议动作完全相反。}原文建议 S480 向 S477 学习、每周多投入
4.1 小时；本文的模型显示这样做会使上行强度比从 0.703 降到 0.669，跃层概率
从 0.642 降到 0.617。相反，只花 1.5 小时提升兴趣与课外拓展，可把上行强度
比提高到 0.954，跃层概率提高到 0.781。这一分歧不是计算误差，而是两种评价
逻辑的必然结果：DEA 只看“投入—产出比”，把人当作黑箱；随机过程把“投入如
何转化为状态转移”写成参数，能区分“有效的投入”与“稀释比值的投入”。

\subsubsection{互补而非替代}

如果必须选一个，本文的答案是“按问题选工具”，而不是“择一而废”：

\begin{itemize}[leftmargin=2em,itemsep=2pt,topsep=2pt]
\item 需要\textbf{横向比较}（谁更接近前沿、哪一类学生整体偏低）时，DEA 更
  合适：它无需分布假设，计算稳定，结果直观。
\item 需要\textbf{纵向判断}（多久能改善、改善的概率多大、会不会回落）时，
  随机过程更合适：它把时间写进模型，输出的是可执行的日程而非排名。
\item 需要\textbf{决策}（有限学时投到哪一项）时，只有把两者结合才完整：先
  用 DEA 定位“谁需要干预、差在哪一维”，再用随机过程的 $\delta$ 筛选“哪一
  维的干预真的提高上行强度”，最后用优化给出“投多少”。
\end{itemize}

因此，本文并不主张“随机过程优于 DEA”，而主张\textbf{DEA 负责“定位”，随
机过程负责“定时”，优化负责“定策”}。原文的问题不在于用了 DEA，而在于把
DEA 的输出当成了终点。

% =============================================================================
\section{结论}

本文以一篇发表于 \textit{Expert Systems with Applications} 的学习策略推荐
论文为对象，完成了从解析、复现到方法替换的完整闭环。主要结论如下。

\textbf{（1）原文的方法链可以被完整拆解，且在结构上可以被复现。}原文的
“复合量化—聚类—类内 DEA—分层—三准则—参考路径”五环链条，每一环都可以写
成明确的数学形式（式\eqref{eq:quant}—\eqref{eq:H}）与算法（算法 1 与算法
2）。在锚点约束下重建的 519 份样本上，本文复
现出同构的五类学生结构（第 5 类人数最多、效率最高）、同向的分层结构与同向
的参考路径；S480 的类内效率精确命中 0.6771，单步增益 26.31\% 与原文
25.38\% 相差 0.93 个百分点，累计增益 47.69\% 与 43.97\% 相差 3.72 个百分点。
\textbf{结构性结论复现成立，数值性结论在锚点约束下成立。}

\textbf{（2）原文存在三处结构性局限与一处内部不一致。}三处局限分别是：评
价是截面的、缺少时间维度；层级依赖参照集、不可跨群体迁移；推荐规则本质是
模仿同伴、不是最优控制。一处不一致是：文字定义 $H=1/\Omega+P$，但表 3 中
的 $H$ 值无法由同一行的 $(\Omega,P)$ 生成。这四点不是吹毛求疵，而是“把
DEA 的输出直接当作建议”这一做法的系统性风险。

\textbf{（3）用随机过程替换静态评价是可行的，且能补上原文缺失的三类信
息。}把五级层级视为状态空间后，带协变量的生灭过程可以估计出上行—下行强
度比 $\rho_i$；用 $\rho_i$ 可以算出任意时长的转移概率、首达时间的解析解与
期望效率。S480 的 16 周跃层概率为 0.642、期望效率提升 25.53\%，与原文的
25.38\% 相互印证；首达最高层的期望时间为 14.5 周，累计 43.97\% 的增益在
连续时间下需要 53.7 周。这些数字在原文中完全缺失。

\textbf{（4）“下行锁定”是群体性现象。}98.1\% 的 $\rho_i<1$，意味着在给定
行为结构的条件下，绝大多数学生的效率层级净漂移向下；模型平稳分布把最大质
量放在第 5 层，与观测到的 45.1\% 相呼应。这一结论的政策含义是：\textbf{不
改变行为结构的“自然等待”只会让人停在低效层，干预必须改变 $\rho$ 本身}。
鞅与停时定理把它进一步落成概率：从第 4 层出发，先跌回第 5 层的概率是
86\%，先爬到第 1 层的概率只有 14\%；从第 3 层出发是 67\% 对 33\%。这组数字
只依赖 $\rho$ 的相对大小，不依赖外部锚点，是截面数据能支撑的最强陈述。

\textbf{（5）“加大投入”不等于“提高效率”，这是本文与原文最重要的分歧。}
在 DEA 的效率口径下，增加投入而产出不变会稀释比值，因此 $\delta_{\text{复
习}}$、$\delta_{\text{作息}}$ 为负；相反，“兴趣与课外拓展”的 $\delta$ 为
正，是唯一能显著提高上行强度的杠杆。据此，模型给出的最优策略是“1.5 小时
兴趣拓展”，而非原文建议的“每周 4.1 小时模仿高效同伴”。后者在模型下的跃
层概率反而低于不干预。

\textbf{（6）两种范式应当互补使用。}DEA 擅长“定位”（谁在前沿、差在哪
维），随机过程擅长“定时”（多久、多大把握），预算约束优化擅长“定策”
（投哪项、投多少）。把三者串起来，才构成一条从数据到建议的完整链条。

\textbf{（7）替换方法的每一步都有可核对的数字。}五种工具各对应一处校验：一致化（Poisson 计数 +
跳链）重建 $\bP(16)$ 与矩阵指数的误差 $3.8\times10^{-15}$；鞅停时定理论证的
$\E[T_1]=m(S_0)$ 与三对角方程组、Gillespie 模拟相差 $0.15\%$；Itô 公式给出
的平稳标准差与类内实测逐类相符；反射原理对纯 Brown 运动给出精确的 $1.5\%$。
反过来，把反射原理照搬到 OU 上会把 $0.105$ 抬到 $0.353$——这一“同族不同律”
的边界必须写明，否则公式就会冒充结论。

本文的局限需要明确列出。\textbf{其一，数据是标定重建的}，因此本文对原文
的“复现”是结构同构意义上的复现，不能替代对原始问卷的复现；所有数值结论
都应理解为“在给定生成机制下的推论”。\textbf{其二，转移强度只被部分识
别}，截面似然只能定出 $\rho$ 的相对大小与符号，时间尺度依赖一个外部锚点，
个体层面的层级预测准确率仅 45.7\%。\textbf{其三，模型依赖较强的假设}：
时间齐次、马尔可夫、层内同质、投入与转化能力可分离。若真实数据中投入与能
力高度相关，$\delta$ 的符号会改变，第（5）条结论需要重新检验。\textbf{其
四，未纳入教师、课程与同伴效应}，而这些因素在教育情境中往往比个体行为更
重要。

后续工作有三个方向。第一，收集三到四个学期的纵向数据（教务系统与辅导员记
录即可），用 Kalbfleisch--Lawless 型面板似然替代截面似然，从而识别个体层
面的转移强度，并检验时间齐次假设。第二，引入隐马尔可夫或多状态模型处理
“层级观测有误差”的问题，把测量误差与真实转移分开。第三，把式
\eqref{eq:opt} 从“一次性动作选择”扩展为“学期内的分阶段预算分配”，用最优
停时或随机控制求解，使建议从“这学期做什么”细化到“第几周做什么、做多
久”。这三项都不需要推翻本文的框架，只需要补上数据的维度。




